\section{Linear Mappings of Finite Rank}

Let A be a ring, and let E and F be A-modules. We denote by \(\operatorname{Hom}_{\mathrm{A}}^{\mathrm{f}}(\mathrm{E},\mathrm{F})\) the set of linear mappings from E to F whose image is contained in a finitely generated submodule of F. It is a subgroup of \(\operatorname{Hom}_{\mathrm{A}}(\mathrm{E},\mathrm{F})\) . When A is a field, it is the set of linear mappings from E to F of finite rank (II, §7, No.~4, p.~298, Definition 3). Set \(\operatorname{End}_{\mathrm{A}}^{\mathrm{f}}(\mathrm{E}) = \operatorname{Hom}_{\mathrm{A}}^{\mathrm{f}}(\mathrm{E},\mathrm{E})\) .

Let E, F, G be A-modules, and let \(u : E \to F\) and \(v : F \to G\) be A-linear mappings. If \(u \in \operatorname{Hom}_{\mathrm{A}}^{\mathrm{f}}(\mathrm{E}, \mathrm{F})\) or \(v \in \operatorname{Hom}_{\mathrm{A}}^{\mathrm{f}}(\mathrm{F}, \mathrm{G})\) , then \(v \circ u\) belongs to \(\operatorname{Hom}_{\mathrm{A}}^{\mathrm{f}}(\mathrm{E}, \mathrm{G})\) .

Denote by \(\theta\) the canonical group homomorphism from \(\mathrm{E}^* \otimes_{\mathrm{A}} \mathrm{F}\) to \(\mathrm{Hom}_{\mathrm{A}}(\mathrm{E}, \mathrm{F})\) (II, §4, No.~2, p.~271); it sends an element \(x^* \otimes y\) of \(\mathrm{E}^* \otimes_{\mathrm{A}} \mathrm{F}\) to the mapping \(x \mapsto \langle x, x^* \rangle y\) from \(\mathrm{E}\) to \(\mathrm{F}\) .

Lemma 1. --- Suppose that the A-module F is projective. The group homomorphism \(\theta\) is injective, and its image is \(\mathrm{Hom}_{\mathrm{A}}^{\mathrm{f}}(\mathrm{E},\mathrm{F})\) .

By loc. cit., Corollary, the homomorphism \(\theta\) is injective. Its image is contained in \(\operatorname{Hom}_{\mathrm{A}}^{\mathrm{f}}(\mathrm{E},\mathrm{F})\) . Let \(u \in \operatorname{Hom}_{\mathrm{A}}^{\mathrm{f}}(\mathrm{E},\mathrm{F})\) ; let us prove that u belongs to the image of \(\theta\) .

First, suppose that the A-module F is free. Let \((f_{i})_{i\in\mathrm{I}}\) be a basis of F, and let \((f_{i}^{*})_{i\in\mathrm{I}}\) be the basis of \(F^{*}\) dual to \((f_{i})_{i\in\mathrm{I}}\) . Let J be a finite subset of I such that the image of u is contained in the submodule of F generated by the \(f_{j}\) for \(j\in J\) . We have

\[
u (x) = \sum_ {j \in \mathrm{J}} \langle u (x), f _ {j} ^ {*} \rangle f _ {j} = \sum_ {j \in \mathrm{J}} \langle x, ^ {t} u (f _ {j} ^ {*}) \rangle f _ {j}\tag{1}
\]

for every \(x\in \mathrm{E}\) , and therefore

\[
u = \theta \left(\sum_ {j \in \mathrm{J}} ^ {t} u \left(f _ {j} ^ {*}\right) \otimes f _ {j}\right).\tag{2}
\]

In the general case, there exist a free A-module L and homomorphisms \(i: F \to L\) and \(p: L \to F\) such that \(p \circ i = 1_{F}\) . The homomorphism \(i \circ u\) belongs to \(\operatorname{Hom}_{\mathrm{A}}^{\mathrm{f}}(\mathrm{E}, \mathrm{L})\) . By the above, there exist a finite set J and, for every \(j \in J\) , elements \(x_{j}^{*}\) of \(E^{*}\) and \(y_{j}\) of L such that we have

\[
i (u (x)) = \sum_ {j \in \mathrm{J}} \langle x, x _ {j} ^ {*} \rangle y _ {j}
\]

for every \(x \in \mathrm{E}\) . We then have

\[
u (x) = p \left(i (u (x))\right) = \sum_ {j \in J} \langle x, x _ {j} ^ {*} \rangle p \left(y _ {j}\right)
\]

for every \(x \in \mathrm{E}\) , and therefore \(u = \theta\big(\sum_{j \in \mathrm{J}} x_j^* \otimes p(y_j)\big)\) .

